% SPDX-FileCopyrightText: 2026 Will Cook
% SPDX-License-Identifier: CC-BY-4.0
%
% Margin links to the formal evidence for each result.
%
% Each paper inputs its own generated file paper/evidence/<paper>.tex, written by
% scripts/paper_evidence.py from the coverage ledger and the Comparator replay
% receipts.  That file holds one line per result,
%
%   \DeclareResultEvidence{<label>}{<mark>}{<Lean URL>}{<Comparator URL>}
%
% and nothing else.  When the result's \label is set, its marks are placed in
% the right margin beside it: "Lean" (or "Lean" with a dagger when the Lean
% proof assumes a named input) and, when the result was compared, "Comparator".
% No text is added to the page.  A declared result whose label never occurs is
% an error, so a mark cannot silently go missing; so is a label declared twice.
% Do not edit the per-paper files by hand.

\usepackage{marginnote}
\renewcommand*{\marginfont}{\footnotesize\sffamily\upshape}
\setlength{\marginparwidth}{56pt}
\setlength{\marginparsep}{14pt}

\makeatletter
\newcommand*{\ev@all}{}
\newcommand{\DeclareResultEvidence}[4]{%
  \ifcsname ev@map@#1\endcsname
    \PackageError{paper-evidence}{Evidence declared twice for label #1}%
      {Regenerate the paper's evidence file with scripts/paper_evidence.py.}%
  \else
    \expandafter\gdef\csname ev@map@#1\endcsname{\ev@marks{#1}{#2}{#3}{#4}}%
    \g@addto@macro\ev@all{\do{#1}}%
  \fi}
\newcommand{\ev@marks}[4]{%
  \ifcsname ev@placed@#1\endcsname\else
    \expandafter\gdef\csname ev@placed@#1\endcsname{}%
    \if\relax\detokenize{#3}\relax
      \PackageError{paper-evidence}{No Lean target for #1}%
        {Regenerate the paper's evidence file with scripts/paper_evidence.py.}%
    \else
      \ev@place{\raggedright\href{#3}{#2}%
        \if\relax\detokenize{#4}\relax\else\\[1pt]\href{#4}{Comparator}\fi}%
    \fi
  \fi}
% \marginnote sits level with the line where the label is set, which is the
% result's heading line: every label of a result is set on its \begin line or at
% the start of the line after it.  scripts/check_paper_evidence_pdfs.py checks
% this on the rendered pages.
\newcommand{\ev@place}[1]{\marginnote{#1}}
\newcommand{\ev@hook}[1]{%
  \ifcsname ev@map@#1\endcsname\csname ev@map@#1\endcsname\fi}
\AtBeginDocument{%
  \let\ev@label\label
  \renewcommand{\label}[1]{\ev@label{#1}\ev@hook{#1}}%
  \ifdefined\compatlabel
    \let\ev@compatlabel\compatlabel
    \renewcommand{\compatlabel}[1]{\ev@compatlabel{#1}\ev@hook{#1}}%
  \fi}
\AtEndDocument{%
  \def\do#1{\ifcsname ev@placed@#1\endcsname\else
    \PackageError{paper-evidence}{The evidence mark for #1 was never placed}%
      {Its label does not occur in this paper; regenerate the evidence file.}\fi}%
  \ev@all}
\makeatother

% Once per paper: what the marks mean.  #1 is the URL of the paper's evidence
% record.
\newcommand{\evidenceparagraph}[1]{%
  \par\smallskip\noindent\emph{Formal proofs.}
  A result with a kernel-checked Lean proof carries a mark in the margin.
  \emph{Lean} opens the proof: the declaration itself when one declaration
  states the whole result, otherwise the list of declarations that together
  state it.  \emph{Comparator} opens the record of an independent check, in
  which the same statement, written again from Mathlib alone in a separate
  repository, was compared with our proof by Lean's Comparator tool, allowing
  only the three standard axioms.  A dagger on the Lean mark means that the
  Lean proof assumes an input named just below the result.  A result without a
  mark has no Lean proof of its whole statement; what is checked is said below
  it.  The \href{#1}{evidence record} gives every declaration, version and
  check.  These checks show that the stated propositions are proved; whether
  they are the right propositions is a question the reader can settle by
  comparing them with the text.}

% A plain remark after a result, for what its mark alone cannot say: the input
% a conditional Lean proof assumes, or what is and is not checked for a result
% without a mark.
\newcommand{\evidenceremark}[1]{%
  \par\nobreak\vspace{-0.3\baselineskip}%
  \noindent{\small #1\par}\vspace{0.15\baselineskip}}
